\section{Conclusion}
\label{sec:conclusion}

This survey organized COF photocatalytic hydrogen evolution by performance factor rather than
material family. The factor chain --- linkage chemistry (L1), band engineering (L2), exciton
and order management (L3), interfacial catalysis (L4), synthesis methodology (L5) --- maps each
design lever to the pipeline stage it governs, with crystallinity and charge extraction
established as the prime factors by the field's clearest causal
benchmark~\cite{ghosh2020identification}. Along that chain, the $\beta$-ketoenamine tautomer
lock stands out as the load-bearing chemical innovation for aqueous operation: a reversible,
error-correcting condensation that ends in an irreversible, hydrolysis-proof
sink~\cite{kandambeth2012construction}.

The survey's second half converted that evidence into practice. The six-entry gap register
shows where the literature stops --- most actionably, the green and scalable synthesis routes
whose HER consequences remain unmeasured --- and the Route~C plan turns that specific gap into
an executable protocol: a solvothermal baseline with two quantified optimizations (microwave:
1~h at 100~$^\circ$C with a $\sim$4.8$\times$ porosity gain; green-solvent flow: 30$\times$
space-time yield at $-$89\% specific energy), a thermodynamic audit that refuses to invent
numbers, and a mandatory safety annex~\cite{grenu2020microwave,xu2026structural}.

The reader we designed for should leave with two assets: a mental model of which factor limits
which performance regime, and a worked template for converting survey evidence into a synthesis
decision --- the survey-to-synthesis bridge that the COF hydrogen-evolution field, now rich in
materials and mechanisms but poor in standardized decisions, most needs.
